\section{A silent center and a signed rim}\label{sec:coaxial}
The simplest test separates a point observation from a measurement around a loop. A steady source circle and a coaxial receiving circle can have a zero vector observation at the receiving circle's center while the oriented rim integral is strictly signed. This is the first result to explain. We then ask how a native state can prepare the source and whether the signed signal survives at a receiver with the model's actual finite radius.

The preceding GM01--GM03 paper imposed Newtonian and magnetostatic field laws through an explicit source dictionary and studied interaction-driven feedback~\cite{GM123}. The present paper follows a different chain: the adopted C6 exterior kernel, a prepared carrier, a proposed native-to-mean impulse, and a passive contour readout~\cite{GM04,GM05,GM06}. A line circle has finite-radius corrections; it is not the radial magnetization profile for which the earlier paper obtained an exact external dipole interaction.

\subsection{The adopted observation and its orientation}
Write a source circle and its oriented unit tangent as
\begin{equation}\label{eq:sourcecircle}
\gamma(\phi)=Y+a(\cos\phi\,e_1+\sin\phi\,e_2),\quad
\widehat t(\phi)=-\sin\phi\,e_1+\cos\phi\,e_2,\quad
\nu=e_1\times e_2,
\end{equation}
where \(e_1,e_2\) are orthonormal, \(a>0\), and \(L=2\pi a\).
C6 adopts a straight delay, tangent encoding and inverse-distance kernel~\cite{PAIRII,INT}:
\begin{align}
U(P,t)&=\frac{g_o}{4\pi}\int_0^L
\frac{p(s,t-d_P(s)/c_e)}{d_P(s)}\dd s,\notag\\
V(P,t)&=\frac{g_o}{4\pi}\int_0^L
\frac{r_c(s,t-d_P(s)/c_e)\widehat t(s)}{d_P(s)}\dd s,
\quad d_P(s)=|P-\gamma(s)|. \label{eq:C6}
\end{align}
Here \(g_o,c_e>0\). A registered point probe returns \((U,B^TV)\) for a known orthonormal frame B. The tangent-encoded \(r_c\widehat t\) is not assigned material electric current, and C6 is a passive observation without loading.

Prepare the steady source \(p=0,r_c=j\), with the corresponding constant history. Then
\begin{equation}\label{eq:staticV}
V(P)=\frac{g_oj}{4\pi}\oint_{\Gamma}\frac{\dd\gamma}{|P-\gamma|},
\qquad H(P):=\nabla_P\times V(P).
\end{equation}
The symbol H denotes this mathematical curl. Identifying it with a measured magnetic field would require calibration beyond the adopted kernel.

\begin{figure}[tbp]
\centering\begin{tikzpicture}[x=1cm,y=1cm,>=Latex,font=\small]
\fill[teal!7] (6,0) ellipse (.58 and 1.45);
\draw[ink,very thick] (0,0) ellipse (.58 and 1.45);
\draw[teal,very thick] (6,0) ellipse (.58 and 1.45);
\draw[dashed,gray] (-.9,0)--(7.2,0);
\draw[->,ink] (0,1.45) arc[start angle=90,end angle=40,x radius=.58,y radius=1.45];
\draw[->,teal] (6,1.45) arc[start angle=90,end angle=40,x radius=.58,y radius=1.45];
\draw[<->] (.2,1.95)--node[fill=white] {$z\ne0$}(5.8,1.95);
\fill[ink] (0,0) circle(1.6pt);
\fill[teal] (6,0) circle(1.6pt);
\node[anchor=east] at(-.15,.12) {$Y$};
\node[anchor=west] at(6.12,.12) {$X$};
\node[ink] at(0,2.55) {steady source};
\node[teal] at(6,2.55) {receiving rim};
\draw[->,ink] (.25,-.55)--(1.6,-.55) node[right] {$\nu$};
\node[ink,align=center] at(1.5,-1.3) {$r=j,\ p=0$\\radius $a$};
\node[teal,align=center] at(4.25,.85) {disk $D_b$\\radius $b$};
\draw[teal] (4.95,.55)--(5.65,.35);
\fill[rust] (6,-2.05) circle(2pt);
\draw[densely dotted] (6,-.1)--(6,-1.95);
\node[anchor=west,rust,align=left] at(6.15,-2.05) {$O_R$\\body-center probe};
\draw[<->,rust] (7.75,0)--node[right] {$wK_c$}(7.75,-2.05);
\draw[->,gray] (-1.6,-1.6)--(-1.6,-.55) node[above] {$e_1$};
\draw[->,gray] (-1.6,-1.6)--(-.7,-1.6) node[right] {$\nu$};
\node[align=center] at(2.7,-2.6) {$V(X)=0$ at the circle center\\$S_b\ne0$ for $j\ne0$};
\end{tikzpicture}

\caption{Coaxial test, shown schematically in an oblique projection; not to scale. Increasing source and receiver angles use the same oriented plane, with \(\nu=e_1\times e_2\) pointing along the common axis. The receiving disk center X is on that axis. Its body-center probe \(O_R=X-wK_ce_1\) is displaced and can already read a nonzero V. The disk is an integration device, not a material panel.}\label{fig:coaxial}
\end{figure}

\subsection{The finite receiving circle}
Place the receiving contour at
\begin{equation}\label{eq:coaxreceiver}
X=Y+z\nu,\qquad q(\theta)=X+b(\cos\theta\,e_1+\sin\theta\,e_2),
\quad t(\theta)=-\sin\theta\,e_1+\cos\theta\,e_2,
\end{equation}
with \(b>0,z\ne0\). Increasing theta fixes the measured orientation. Its area-normalized signed reading is
\begin{equation}\label{eq:sensor}
S_b(X,\nu)=\frac1{\pi b^2}\oint_{\partial D_b}V\cdot\dd\ell
=\frac1{\pi b}\int_0^{2\pi}V(q(\theta))\cdot t(\theta)\dd\theta.
\end{equation}
At X every source point has distance \(\sqrt{a^2+z^2}\), and \(\oint\dd\gamma=0\); consequently \(U(X)=V(X)=0\), for both an empty source and every constant j.

\begin{theorem}[Finite coaxial sign]\label{thm:coaxial}
With the orientation above, \(S_b\) has the sign of j for every \(a,b>0,z\ne0\). It vanishes at \(j=0\). Reversing j or the receiving traversal reverses the reading.
\end{theorem}
\begin{proof}
The double line integral has tangent pairing \(ab\cos(\phi-\theta)\) and squared distance
\(a^2+b^2+z^2-2ab\cos(\phi-\theta)\). Rotational symmetry removes the theta integral, giving
\begin{equation}\label{eq:coaxexact}
S_b=\frac{g_oj a}{2\pi b}\int_0^{2\pi}
\frac{\cos\phi}{\sqrt{C-q_0\cos\phi}}\dd\phi,\qquad
C=a^2+b^2+z^2,\quad q_0=2ab.
\end{equation}
Since \(C-q_0=(a-b)^2+z^2>0\), pairing opposite cosines rewrites the integral as
\begin{equation}\label{eq:positivepair}
\int_{-\pi/2}^{\pi/2}\cos\phi
\left[(C-q_0\cos\phi)^{-1/2}-(C+q_0\cos\phi)^{-1/2}\right]\dd\phi>0.
\end{equation}
The integrand is positive in the interval interior. Linearity in j and orientation of the line integral prove the remaining statements.
\end{proof}

\begin{corollary}[One rim sample in this control]\label{cor:onesample}
For this steady aligned coaxial source only, every rim point gives
\begin{equation}\label{eq:onesample}
S_b=\frac2b\,V(q(\theta))\cdot t(\theta).
\end{equation}
\end{corollary}
\begin{proof}
Rotation about the common axis carries both \(V(q(\theta))\) and \(t(\theta)\) by the same rotation. Their dot product is constant; substitute it into \eqref{eq:sensor}.
\end{proof}
One ideal signed tangential sample thus suffices in this symmetric control. This is not a universal replacement for a contour or the later three-source sampling certificate.

Integrating the derivative of \(x^{-1/2}\) in the paired expression gives, for \(j>0\),
\begin{equation}\label{eq:coaxbounds}
\frac{g_oj a^2}{2[z^2+(a+b)^2]^{3/2}}
\le S_b\le
\frac{g_oj a^2}{2[z^2+(a-b)^2]^{3/2}}.
\end{equation}
The positivity is exact at finite radii. The bounds are conservative.

\subsection{A circle on an actual copy}
The established geometry supplies~\cite{PAIR}
\begin{equation}\label{eq:geometryconstants}
K_c=\frac{\sqrt3}{12}+\frac{\sqrt6}{4},\qquad
R_\Gamma=\frac{\sqrt7K_c}{2(1+\sqrt3/2)},\qquad
a=b=wR_\Gamma,\quad R_\Gamma=0.53645209718\ldots,
\end{equation}
where w is common host width. In the published orientation, \(Y=O+wK_c\ey\), \(e_1=\ey,e_2=\ez,\nu=\ex\). Translate an unrotated same-width copy by \(z\ex\). Its circle center is \(X=Y+z\ex\), but its body center is \(O_R=O+z\ex=X-wK_c\ey\).

This is an explicit rigid test placement, preserving the established circle and full-copy geometry; it is not one of the published H/R array fixtures. The zero reading is at the \emph{circle center}, not generally at the existing body-center probe. At the benchmark \(z=4w\), enclosing balls of radius \(w(K_c+R_\Gamma)<2w\) separate the two full copies and all their circles. The all-\(z\ne0\) integral theorem alone does not assert full-body noncollision at small z.
